\section{Sandeep}

A critical question in understanding the labor market and informing policy recommendations , is whether the observed differences in earnings in the formal and informal employment is the result of market segmentation or the incentives within competitive market framework.  In other words, the question is whether individuals opt for informal employment by choice or are pushed into it due to the barriers of entry into formal employment. Although the question is particularly important in developing economies, where the phenomenon of informality is widespread; no consensus has been achieved on the driving mechanism of informality due to the complex structure of informal employment.

The traditional dual market theory, starting with (Lewis, 1954) conceive informal employment as backward, marginal and disadvantaged sector where employees are paid wages above market-clearing prices and serves as a last resort to escape unemployment (Harris and Todaro, 1970, FIelds, 1990, 2005a, Stighlitz, 1976). According to this approach, workers in the informal sector is predominantly involuntarily engagement of workers and earn less than identical workers in the formal sector. This segmentation may result from informational or geographic barriers to entry, or from sector-specific variations in employment and wage regulations (Dikens and lang, 1992). Firms may offer wages above market clearing prices to increase profit by increasing retention, reducing turnover costs, eventually maximizing productivity of employees     


     

   













In Figure 1, we present estimates of the corresponding quantities from the quantile decompositions across the overall distribution of wage gap attributable to differences in the average characteristics of the two groups (composition effect) and differences in returns to those characteristics (structure effect). These estimates are detailed for the overall service sector as well as its stratified segments—upper tier, middle tier, and lower tier—across two distinct years, 2008 and 2018. The stratified segment reveal heterogeneity in the relative importances of covariates in accounting for disparities across the distribution, which the unidimensional classical method cannot capture. The results reveal that the wage gap is maximum below the median ($\tau$ 0.25 - 0.5), whereas it deteriorates at an increasing rate above.  



 